\chapter{Easy Instrumentation with \eacsl}

Due to numerous instrumentation and compile-time options generating monitored
executables using the \eacsl plug-in can be tedious. To simplify this task we
developed \eacslgcc\ -- a convenience wrapper around \framac and \gcc allowing
to instrument programs with \eacsl and compile the instrumented code using a
single command and a very few parameters.

In this section we describe \eacslgcc and provide several examples of its use.

%%%%%%%%%%%%%%%%%%%%%%%%%%%%%%%%%%%%%%%%%%%%%%%%%%%%%%%%%%%%%%%%%%%%%%%%%%%%%%%
%%%%%%%%%%%%%%%%%%%%%%%%%%%%%%%%%%%%%%%%%%%%%%%%%%%%%%%%%%%%%%%%%%%%%%%%%%%%%%%
%%%%%%%%%%%%%%%%%%%%%%%%%%%%%%%%%%%%%%%%%%%%%%%%%%%%%%%%%%%%%%%%%%%%%%%%%%%%%%%

\section{Basic Use}

\eacslgcc instruments C files provided at input using the \eacsl plug-in, and
optionally compiles the generated code and the original sources using a \C
compiler (\gcc by default).

The default behaviour of \eacslgcc is to instrument an annotated C program
with runtime assertions via a run of \framac.

\begin{shell}
 \$ e-acsl-gcc.sh -ogen.c foo.c
\end{shell}

The above command instruments \texttt{foo.c} with runtime assertions and
outputs the generated code to \texttt{gen.c}. Unless the name of the output
file is specified via the \texttt{-o} option (i.e., \texttt{-ogen.c} in the
above example), the generated code is placed to a file named
\texttt{a.out.frama.c}.

Compilation of the original and the instrumented sources
is enabled using the \texttt{-c} option.  For instance, the command

\begin{shell}
 \$ e-acsl-gcc.sh -c -ogen.c foo.c
\end{shell}

instruments the code in \texttt{foo.c}, outputs it to \texttt{gen.c} and
produces 2 executables: \texttt{a.out} and \texttt{a.out.e-acsl}, where
\texttt{a.out} is an executable compiled from \texttt{foo.c}, while
\texttt{a.out.e-acsl} is an executable compiled from \texttt{gen.c} generated
by \eacsl.

Named executables can be created using the \texttt{-O} option.  This is such
that a value supplied to the \texttt{-O} option is used to name the executable
generated from the original program and the executable generated from the
\eacsl-instrumented code is suffixed \texttt{.e-acsl}.

For example, command
\begin{shell}
 \$ e-acsl-gcc.sh -c -ogen.c -Ogen foo.c
\end{shell}
names the executables generated from \texttt{foo.c} and \texttt{gen.c}:
\texttt{gen} and \texttt{gen.e-acsl} respectively.

The \eacslgcc \texttt{-C} option allows to skip the instrumentation step and
compile a given program as if it were already instrumented by the \eacsl
plug-in.  This option is useful if one makes changes to an already
instrumented source file by hand and only wants to recompile it.

\begin{shell}
 \$ e-acsl-gcc.sh -C -Ogen gen.c
\end{shell}

In the above example the run of \eacslgcc produces a single executable named
\texttt{gen.e-acsl}. This is because \texttt{gen.c} is considered to be
instrumented code and thus no original program is provided.

%%%%%%%%%%%%%%%%%%%%%%%%%%%%%%%%%%%%%%%%%%%%%%%%%%%%%%%%%%%%%%%%%%%%%%%%%%%%%%%
%%%%%%%%%%%%%%%%%%%%%%%%%%%%%%%%%%%%%%%%%%%%%%%%%%%%%%%%%%%%%%%%%%%%%%%%%%%%%%%
%%%%%%%%%%%%%%%%%%%%%%%%%%%%%%%%%%%%%%%%%%%%%%%%%%%%%%%%%%%%%%%%%%%%%%%%%%%%%%%

\section{Documentation}

A list of \eacslgcc options can be retrieved using the following command:

\begin{shell}
 \$ e-acsl-gcc.sh -h
 Usage: e-acsl-gcc.sh [options] files
 Options:
   -h         show this help page
   -c         compile instrumented code
   -l         pass additional options to the linker
   -e         pass additional options to the pre-preprocessor
   -E         pass additional arguments to the Frama-C pre-processor
   -p         output the generated code with rich formatting to STDOUT
   -o <file>  output the generated code to <file> [a.out.frama.c]
   -O <file>  output the generated executables to <file> [a.out, a.out.e-acsl]
   -M         maximize memory-related instrumentation
   -g         always use GMP integers instead of C integral types
   -q         suppress any output except for errors and warnings
   -s <file>  redirect all output to <file>
   -P         compile executable without debug features
   -I <file>  specify Frama-C executable [frama-c]
   -G <file>  specify GCC executable [gcc]
\end{shell}

Note that the \texttt{-h} option lists only the most basic options.
For full up-to-date documentation of \eacslgcc see its man page:

\begin{shell}
    \$ man e-acsl-gcc.sh
\end{shell}

%%%%%%%%%%%%%%%%%%%%%%%%%%%%%%%%%%%%%%%%%%%%%%%%%%%%%%%%%%%%%%%%%%%%%%%%%%%%%%%
%%%%%%%%%%%%%%%%%%%%%%%%%%%%%%%%%%%%%%%%%%%%%%%%%%%%%%%%%%%%%%%%%%%%%%%%%%%%%%%
%%%%%%%%%%%%%%%%%%%%%%%%%%%%%%%%%%%%%%%%%%%%%%%%%%%%%%%%%%%%%%%%%%%%%%%%%%%%%%%

\section{Customising \framac and \gcc Invocations}

Runs of \framac and \gcc issued by \eacslgcc can be customized by using
additional flags.

Additional \framac flags can be passed using the \texttt{-F} option, while
\texttt{-l} and \texttt{-e} options allow for passing additional pre-processor
and linker flags during \gcc invocations.  For example, command

\begin{shell}
 \$ e-acsl-gcc.sh -c -F"-unicode" -e"-I/bar -I/baz" foo.c
\end{shell}

yields an instrumented program \texttt{a.out.frama.c} where ACSL formulas
are output using the UTF-8 character set. Further, during the compilation of
\texttt{a.out.frama.c} \gcc adds directories \texttt{/bar} and
\texttt{/baz} to the list of directories to be searched for header files.

It is worth mentioning that \eacslgcc implements several convenience flags for
setting some of the \framac options. For instance, \texttt{-E} can be used to
pass additional arguments to the \framac pre-processor and \texttt{-M}
maximizes memory-related instrumentation.

%%%%%%%%%%%%%%%%%%%%%%%%%%%%%%%%%%%%%%%%%%%%%%%%%%%%%%%%%%%%%%%%%%%%%%%%%%%%%%%
%%%%%%%%%%%%%%%%%%%%%%%%%%%%%%%%%%%%%%%%%%%%%%%%%%%%%%%%%%%%%%%%%%%%%%%%%%%%%%%
%%%%%%%%%%%%%%%%%%%%%%%%%%%%%%%%%%%%%%%%%%%%%%%%%%%%%%%%%%%%%%%%%%%%%%%%%%%%%%%

\section{Verbosity Levels and Output Redirection}

By default \eacslgcc prints the executed \framac and \gcc commands and pipes
their output to the terminal. The verbosity of \framac output can be changed
using \texttt{-v} and \texttt{-d} options used to pass values to
\texttt{-verbose} and \texttt{-debug} flags of \framac respectively.  For
instance, to increase the verbosity of plug-ins but suppress the verbosity of
the \framac kernel while instrumenting \texttt{foo.c} one can use the following
command:

\begin{shell}
 \$ e-acsl-gcc.sh -v5 -F"-kernel-verbose 0" foo.c
\end{shell}

Verbosity of the \eacslgcc output can also be reduced using the \texttt{-q}
option that suppresses any output except errors or warnings issued by \gcc,
\framac or \eacslgcc itself.

The output of \eacslgcc can also be redirected to a specified file using the
\texttt{-s} option.  For example, the following command will redirect all
output during instrumentation and compilation of \texttt{foo.c} to the
file named \texttt{/tmp/log}.

\begin{shell}
 \$ e-acsl-gcc.sh -c -s/tmp/log foo.c
\end{shell}

%%%%%%%%%%%%%%%%%%%%%%%%%%%%%%%%%%%%%%%%%%%%%%%%%%%%%%%%%%%%%%%%%%%%%%%%%%%%%%%
%%%%%%%%%%%%%%%%%%%%%%%%%%%%%%%%%%%%%%%%%%%%%%%%%%%%%%%%%%%%%%%%%%%%%%%%%%%%%%%
%%%%%%%%%%%%%%%%%%%%%%%%%%%%%%%%%%%%%%%%%%%%%%%%%%%%%%%%%%%%%%%%%%%%%%%%%%%%%%%

\section{Executables}

By default \eacslgcc uses the first \texttt{frama-c} executable found in a
system's path to instrument input programs. The name of the \framac executable
can be changed using the \texttt{-I} option. For instance, to instrument a file
named \texttt{foo.c} using the \texttt{/usr/local/bin/frama-c.byte} executable
one can use the following command:

\begin{shell}
 \$ e-acsl-gcc.sh -I/usr/local/bin/frama-c.byte foo.c
\end{shell}

The name of a \gcc executable can be changed similarly using the \texttt{-G}
option.
